\documentclass[11pt]{article}
\usepackage[T1]{fontenc}
\usepackage[margin=1in]{geometry}
\usepackage{amsmath,amssymb}
\usepackage{hyperref}

\title{A Literature-Implied Answer to a Parseval-Frame Class Problem}
\author{Run fa\_banach\_001, lane 2}
\date{2026-07-18}

\begin{document}
\maketitle

\section*{Status}

\textbf{literature\_implied\_answer.} This packet records that a finite
Parseval-frame equivalence-class subproblem stated in arXiv:0906.5606 is
answered by the standard projection/Gramian parametrization of normalized tight
frames, stated explicitly in Holmes--Paulsen \cite{holmespaulsen}.

\section*{Source Problem}

The source is Casazza--Kutyniok--Li--Rozell \cite{cklr}, arXiv:0906.5606.
In Section 4, subsection \emph{Extensions and Related Problems}, under
\emph{Equivalence Classes}, the authors write that the following apparently
simple problem is still unsolved: construct one Parseval frame in each
equivalence class, choosing unitary equivalence as the relation.

\section*{Supporting Literature}

Holmes and Paulsen \cite{holmespaulsen} identify finite normalized tight frames
with rank projections through the analysis operator. In their notation, an
\((n,k)\)-frame has analysis isometry \(V:\mathbb F^k\to\mathbb F^n\), and
\(VV^*\) is the Gramian/correlation projection. They state that there is a
one-to-one correspondence between \(n\times n\) rank \(k\) projections and
type-I equivalence classes of \((n,k)\)-frames. Type-I equivalence is exactly
unitary equivalence of the ordered frame vectors.

\section*{Identification and Representative}

Let \(P\) be a rank \(k\) orthogonal projection on \(\mathbb F^n\), where
\(\mathbb F=\mathbb R\) or \(\mathbb C\). Define
\[
  f_i=P e_i\in \operatorname{ran} P,\qquad i=1,\ldots,n.
\]
Then \(\{f_i\}_{i=1}^n\) is a Parseval frame for \(\operatorname{ran}P\), because
for every \(x\in\operatorname{ran}P\),
\[
  \sum_{i=1}^n |\langle x,Pe_i\rangle|^2
  =
  \sum_{i=1}^n |\langle Px,e_i\rangle|^2
  =
  \sum_{i=1}^n |\langle x,e_i\rangle|^2
  =
  \|x\|^2.
\]
Moreover, the Gram matrix of this frame is
\[
  (\langle f_j,f_i\rangle)_{i,j}
  =
  (\langle Pe_j,Pe_i\rangle)_{i,j}
  =
  P.
\]
Conversely, if \(\{g_i\}_{i=1}^n\) is a Parseval frame for a \(k\)-dimensional
Hilbert space with analysis isometry \(V\), then \(P=VV^*\) is a rank \(k\)
projection and \(Vg_i=Pe_i\). Thus the original frame is unitarily equivalent
to the projection representative \(\{Pe_i\}_{i=1}^n\).

\section*{Scope}

This gives one representative in each finite ordered Parseval-frame
unitary-equivalence class. It does not solve the source's broader fusion-frame
equivalence-class program, nor its separate weighted, chordal-distance, or
sparsity questions.

\begin{thebibliography}{9}

\bibitem{cklr}
P. G. Casazza, G. Kutyniok, S. Li, and C. J. Rozell,
\emph{Fusion frames: existence and construction},
arXiv:0906.5606.

\bibitem{holmespaulsen}
R. B. Holmes and V. I. Paulsen,
\emph{Optimal frames for erasures},
Linear Algebra Appl. 377 (2004), 31--51.
DOI: \href{https://doi.org/10.1016/j.laa.2003.07.012}{10.1016/j.laa.2003.07.012}.

\end{thebibliography}

\end{document}
